\documentclass[11pt,a4paper]{article}

\usepackage[utf8]{inputenc}
\usepackage[T1]{fontenc}
\usepackage{amsmath,amssymb,amsthm,amsfonts}
\usepackage{geometry}
\usepackage{graphicx}
\usepackage{booktabs}
\usepackage{array}
\usepackage{fancyhdr}

% IMPORTANT: load hyperref before styling packages
\usepackage[hypertexnames=false,colorlinks=true,linkcolor=blue,filecolor=magenta,urlcolor=blue,pdftitle={EP001 Physical System Description},pdfauthor={Double-Inverted Pendulum Learning Module}]{hyperref}

\geometry{left=2.5cm,right=2.5cm,top=2.5cm,bottom=2.5cm}

\pagestyle{fancy}
\fancyhf{}
\fancyhead[L]{EP001: Physical System Description}
\fancyhead[R]{Learning Module}
\fancyfoot[C]{\thepage}
\renewcommand{\headrulewidth}{0.4pt}

\newtheorem{theorem}{Theorem}[section]
\newtheorem{definition}{Definition}[section]

\title{\textbf{EP001: Physical System Description}\\\textit{Double-Inverted Pendulum Control Theory}}
\author{Double-Inverted Pendulum Learning Module}
\date{\today}

\begin{document}

\maketitle

\section{Introduction}

\subsection{Overview}
The double inverted pendulum is a classic benchmark in nonlinear control.
\begin{itemize}
    \item It is unstable in open loop.
    \item It has coupled rotational dynamics.
    \item It has one actuator and two rotational states.
\end{itemize}

\subsection{System Characteristics}
The plant is nonlinear and underactuated:
\begin{enumerate}
    \item Non-minimum phase behavior can appear near upright equilibrium.
    \item Full state feedback improves achievable robustness margins.
    \item Disturbances strongly affect transient response.
\end{enumerate}

\section{Physical System Description}

\subsection{System Configuration}
The model consists of a cart, two serial rigid links, and one control force.
\begin{itemize}
    \item Cart mass: $m_c$
    \item Link 1 mass: $m_1$, length: $l_1$, angle: $\theta_1$
    \item Link 2 mass: $m_2$, length: $l_2$, angle: $\theta_2$
\end{itemize}

\subsection{Physical Constraints}
\begin{enumerate}
    \item Cart travel: $x \in [x_{\min}, x_{\max}]$
    \item Force bounds: $u \in [u_{\min},u_{\max}]$
    \item Joint limits near safe mechanical stops if present
\end{enumerate}

\section{Mathematical Model}

\subsection{State Vector Definition}
State vector choice is:
\begin{equation}
    x = [q,\dot q]^T = [x,\theta_1,\theta_2,\dot x,\dot{\theta}_1,\dot{\theta}_2]^T.
\end{equation}

\subsection{Equations of Motion}
A compact form is used:
\begin{equation}
    M(q)\ddot q + C(q,\dot q)\dot q + G(q) = B u + d.
\end{equation}

\begin{definition}
The equation is valid under rigid-link and rigid-joint assumptions.
\end{definition}

\subsection{Inertia Matrix}
The inertia matrix $M(q)$ is symmetric and positive definite.
\begin{equation}
    M(q)=\begin{bmatrix}
    M_{11} & M_{12} & M_{13}\\
    M_{12} & M_{22} & M_{23}\\
    M_{13} & M_{23} & M_{33}
    \end{bmatrix}.
\end{equation}

\subsection{Coriolis and Centrifugal Matrix}
Coriolis terms are configuration and velocity dependent.
\begin{equation}
    C(q,\dot q)\dot q = \begin{bmatrix}c_1\\ c_2\\ c_3\end{bmatrix}.
\end{equation}

\section{Nonlinearity Analysis}

\subsection{Configuration-Dependent Inertia}
The inertia changes with configuration and couples all coordinates.
\begin{enumerate}
    \item Dependence on $\cos\theta_1$, $\cos\theta_2$, and $\cos(\theta_1-\theta_2)$
    \item Determinant and condition number vary with $q$
    \item Ill-conditioning near singular alignment may occur
\end{enumerate}

\subsection{Trigonometric Nonlinearity}
\begin{theorem}
Gravity and coupling create trigonometrical nonlinearities.
\end{theorem}
The dominant terms include $\sin\theta_i$ and $\cos\theta_i$.

\section{System Parameters}

\subsection{Nominal Parameter Values}
\begin{table}[h]
\centering
\caption{Nominal DIP System Parameters}
\label{tab:nominal_parameters}
\begin{tabular}{@{}ll@{}}
\toprule
Parameter & Value \\
\midrule
$m_c$ & 1.0 kg \\
$m_1$ & 0.5 kg \\
$m_2$ & 0.3 kg \\
$l_1$ & 0.4 m \\
$l_2$ & 0.3 m \\
$g$ & 9.81 m/s$^2$ \\
\bottomrule
\end{tabular}
\end{table}

\subsection{Parameter Uncertainty}
\begin{itemize}
    \item Mass and friction values vary with payload.
    \item Length tolerances affect effective lever arms.
    \item Sensing noise biases derivative terms.
\end{itemize}

\section{Linearized Model}

\subsection{Equilibrium Points}
Equilibrium near upright configuration is commonly chosen:
\begin{equation}
    x_e = [0,0,0,0,0,0]^T,
\end{equation}
with other valid equilibria depending on control objective.

\subsection{Linearized State-Space Representation}
After small-angle approximation the local model is:
\begin{equation}
    \dot x = A x + B u + E w.
\end{equation}
\begin{itemize}
    \item Stability of $A$ determines local behavior.
    \item Controller gains depend on performance versus robustness trade-offs.
\end{itemize}

\section{Numerical Properties}

\subsection{State Scaling}
\begin{enumerate}
    \item Angles are typically in radians.
    \item Velocities are in rad/s.
    \item Scaling improves optimization conditioning.
\end{enumerate}

\subsection{Time Scaling}
\begin{enumerate}
    
    \item Choose step size based on highest mode frequency.
    \item Ensure numerical integrator accuracy near impact events.
    \item Use fixed-step simulation for reproducibility.
\end{enumerate}

\section{Summary and Key Takeaways}

\subsection{System Overview}
The double inverted pendulum is a stiff, coupled, underactuated plant.

\subsection{Key Challenges}
\begin{enumerate}
    \item Nonlinearity dominates wide operating ranges.
    \item Strong coupling changes control authority with configuration.
    \item Unmodeled disturbances degrade near-boundary stability.
\end{enumerate}

\subsection{Control Implications}
\begin{theorem}
Control design requires balanced trade-offs between competing objectives.
\end{theorem}
\begin{itemize}
    \item Linear designs simplify synthesis.
    \item Sliding and adaptive methods improve robustness.
    \item Constraint-aware tuning is mandatory in hardware loops.
\end{itemize}

\end{document}
